\section{Conclusion}

In this paper, we proposed GRA-CNN, a novel channel pruning method that leverages Gray Relational Analysis to identify and preserve the most semantically significant filters in deep convolutional neural networks. Unlike traditional magnitude-based or geometric pruning methods that rely on local weight properties or layer-wise statistics, GRA-CNN introduces a global, sequence-level importance metric. By aligning channel activation sequences with the final classification logits, our method effectively captures the contribution of each filter to the network's ultimate decision-making process. This semantic alignment ensures that critical discriminative features are retained even under aggressive compression rates.

Extensive experiments on CIFAR-10, CIFAR-100, and Tiny-ImageNet benchmarks demonstrate the superiority of GRA-CNN over state-of-the-art methods such as L1-Norm, FPGM, and HRank. Specifically, GRA-CNN achieves significantly better accuracy-efficiency trade-offs, particularly in deep architectures like ResNet-56 and ResNet-110. For instance, on ResNet-56, our method maintains 92.00\% accuracy with a 50.5\% reduction in FLOPs, surpassing purely geometric approaches. The robustness of GRA-CNN is further validated by its ability to generalize across different datasets and architectures without requiring complex hyperparameter tuning or extensive retraining.

Future work will focus on extending the GRA framework to other architectures such as MobileNets and Vision Transformers (ViTs), which present different structural challenges. We also plan to explore adaptive reference sequence construction to further refine the semantic alignment process. Additionally, investigating the deployment of GRA-pruned models on resource-constrained edge devices to quantify real-world energy and latency savings remains a promising direction. By bridging the gap between classic system analysis theory and modern neural network compression, GRA-CNN offers a rigorous and effective pathway towards efficient deep learning.
